\documentclass[12pt,a4paper]{article}

% --- PAKETE FÜR DEUTSCHE SPRACHE & SCHRIFTART ---
\usepackage[T1]{fontenc}
\usepackage[ngerman]{babel}
\usepackage{lmodern}
\usepackage{titlesec}

% --- CSQUOTES (Vom Biblatex-Paket empfohlen) ---
\usepackage{csquotes}

% --- MATHE, ENUMERATE, FLOAT & LAYOUT ---
\usepackage{amsmath,amssymb}
\usepackage{geometry}
\geometry{a4paper, margin=2.5cm}
\usepackage{parskip}
\usepackage{enumitem}
\usepackage{float}

% --- TCOLORBOX FÜR HIGHLIGHT- & AUFGABENBOXEN ---
\usepackage[most]{tcolorbox}

\newtcolorbox{taskbox}[1][]{
    colback=blue!5!white,
    colframe=blue!60!black,
    fonttitle=\bfseries,
    title={#1},
    arc=2mm,
    boxrule=0.8pt,
    left=4mm, right=4mm, top=3mm, bottom=3mm
}

\newtcolorbox{solutionbox}[1][Lösung]{
    colback=green!5!white,
    colframe=green!50!black,
    fonttitle=\bfseries,
    title={#1},
    arc=2mm,
    boxrule=0.8pt,
    left=4mm, right=4mm, top=3mm, bottom=3mm
}

\newtcolorbox{emailbox}[1][]{
    colback=red!4!white,
    colframe=red!65!black,
    fonttitle=\bfseries,
    title={#1},
    arc=1.5mm,
    boxrule=0.8pt,
    left=4mm, right=4mm, top=3mm, bottom=3mm
}

% --- FARBEN & VERLINKUNGEN ---
\usepackage{xcolor}
\usepackage[colorlinks=true, linkcolor=blue!60!black, urlcolor=blue!60!black, citecolor=blue!60!black]{hyperref}

% --- CLEVEREF (MUSS NACH HYPERREF GELADEN WERDEN) ---
\usepackage[ngerman]{cleveref}
\crefname{figure}{Abbildung}{Abbildungen}
\Crefname{figure}{Abbildung}{Abbildungen}
\crefname{table}{Tabelle}{Tabellen}
\Crefname{table}{Tabelle}{Tabellen}
\crefname{section}{Abschnitt}{Abschnitte}
\Crefname{section}{Abschnitt}{Abschnitte}

% --- LITERATURVERZEICHNIS MIT BIBLATEX & BIBER ---
\sloppy
\usepackage[backend=biber, style=numeric]{biblatex}
\addbibresource{Skript.bib}

% --- GLOSSAR & ABKÜRZUNGSVERZEICHNIS ---
\usepackage[toc, acronym, nonumberlist]{glossaries}
\makeglossaries

% Custom Glossar-Style für optisch ansprechende Formatierung
\newglossarystyle{mycleanstyle}{%
  \setglossarystyle{long}%
  \renewenvironment{theglossary}%
     {\begin{longtable}{@{}p{5.8cm}@{\hspace{0.6cm}}p{9.6cm}@{}}}%
     {\end{longtable}}%
  \renewcommand*{\glossentry}[2]{%
    \textbf{\glsentryname{##1}} & \glsentrydesc{##1}\\[0.9em]%
  }%
  \renewcommand*{\glsgroupskip}{}%
}
\setglossarystyle{mycleanstyle}

% Inkludieren der separaten Abkürzungs- und Glossardateien
% --- ABKÜRZUNGEN (AKRONYME) ---

\newacronym{mfa}{MFA}{Multi-Faktor-Authentisierung}
\newacronym{nist}{NIST}{National Institute of Standards and Technology}
\newacronym{bsi}{BSI}{Bundesamt für Sicherheit in der Informationstechnik}
\newacronym{smtp}{SMTP}{Simple Mail Transfer Protocol}
\newacronym{url}{URL}{Uniform Resource Locator}
\newacronym{pin}{PIN}{Persönliche Identifikationsnummer}
\newacronym{it}{IT}{Informationstechnik}

% --- GLOSSAREINTRÄGE ---

\newglossaryentry{Multi-Faktor-Authentisierung}{
    name={Multi-Faktor-Authentisierung},
    description={Ein Sicherheitsverfahren, bei dem mehrere unabhängige Faktoren zur Authentifizierung kombiniert werden.}
}

\newglossaryentry{passwort}{
    name={Passwort},
    description={Eine geheime Zeichenkette, die zur Authentisierung dient.}
}

\newglossaryentry{entropie}{
    name={Entropie},
    description={Ein Mass aus der Informationstheorie für die Stärke und Zufälligkeit eines Passworts.}
}

\newglossaryentry{bruteforce}{
    name={Brute-Force-Angriff},
    description={Ein automatisierter Angriff, bei dem Programme alle möglichen Zeichenkombinationen systematisch \mbox{durchprobieren}.}
}

\newglossaryentry{dictionaryattack}{
    name={Dictionary Attack (Wörterbuchangriff)},
    description={Ein gezielter Angriff, bei dem Listen mit \mbox{gebräuchlichen} Wörtern, Namen und bekannten Passwörtern systematisch getestet werden.}
}

\newglossaryentry{datenleck}{
    name={Datenleck},
    description={Ein Vorfall, bei dem vertrauliche Daten unabsichtlich oder absichtlich veröffentlicht oder gestohlen werden.}
}

\newglossaryentry{passwortmanager}{
    name={Passwort-Manager},
    description={Software zur sicheren Verwaltung und Generierung von Passwörtern.}
}

\newglossaryentry{passphrase}{
    name={Passphrase},
    description={Eine längere Zeichenkette, oft bestehend aus mehreren zufälligen Wörtern, die als Passwort verwendet wird.}
}

\newglossaryentry{phishing}{
    name={Phishing},
    description={Eine Form des Social Engineering, bei der Angreifer versuchen, über gefälschte E-Mails, Nachrichten oder Webseiten an vertrauliche Daten zu gelangen.}
}

\newglossaryentry{socialengineering}{
    name={Social Engineering},
    description={Die Manipulation von Personen durch psychologische Tricks, um an vertrauliche Informationen zu gelangen.}
}

\newglossaryentry{spearphishing}{
    name={Spear-Phishing},
    description={Ein gezielter Phishing-Angriff, der auf eine spezifische Person oder bestimmte Mitarbeiter zugeschnitten ist.}
}

\newglossaryentry{quishing}{
    name={Quishing},
    description={Eine Phishing-Methode, bei der manipulierte QR-Codes eingesetzt werden, um Opfer auf gefälschte Webseiten zu leiten.}
}


% --- NEU: ECHTE SUBSUBSUBSECTION DEFINIEREN ---
\titleclass{\subsubsubsection}{straight}[\subsubsection]
\newcounter{subsubsubsection}[subsubsection]
\renewcommand\thesubsubsubsection{\thesubsubsection.\arabic{subsubsubsection}}

\titleformat{\subsubsubsection}
  {\normalfont\normalsize\bfseries}{\thesubsubsubsection}{1em}{}
\titlespacing*{\subsubsubsection}
  {0pt}{3.25ex plus 1ex minus .2ex}{1.5ex plus .2ex}

% Zähler- & TOC-Tiefe auf Ebene 4 anheben
\setcounter{secnumdepth}{4}
\setcounter{tocdepth}{4}

% Eintragen & Einrücken im Inhaltsverzeichnis (TOC)
\makeatletter
\def\toclevel@subsubsubsection{4}
\def\l@subsubsubsection{\@dottedtocline{4}{7.0em}{4.1em}}
\makeatother

\begin{document}

% --- TEIL 1: PASSWORTSICHERHEIT ---
\subsection{Passwortsicherheit}
\label{sec:passwortsicherheit}

\subsubsection{Einleitung}
In fast allen IT-Systemen bilden \glslink{passwort}{Passwörter} die erste Hürde gegen unbefugten Zugriff. In der Praxis zeigt sich jedoch immer wieder, dass schwache, kompromittierte oder mehrfach genutzte \glslink{passwort}{Passwörter} eine der grössten Sicherheitslücken darstellen. Dieses Skript gibt einen Überblick über die mathematischen Grundlagen der \gls{passwort}-Stärke und fasst die aktuellen Empfehlungen zur Absicherung digitaler Identitäten zusammen \cite{nist2024}.

\subsubsection{Grundlegende Begriffe und Angriffsarten}
In der Authentisierung und \gls{passwort}-Sicherheit spielen verschiedene Konzepte und \mbox{Bedrohungen} eine entscheidende Rolle:

Neben dem rein mathematischen \gls{bruteforce} nutzen Angreifer in der Praxis sehr häufig ein \gls{dictionaryattack}, bei dem gezielt Listen mit geläufigen Wörtern, Namen und bekannten Daten \mbox{durchprobiert} werden. Um diesen \mbox{Methoden} zu entgehen, ist eine hohe \gls{entropie} entscheidend. Zusätzlich schützt moderne \mbox{\gls{mfa}} (z.\,B. eine Bestätigung direkt auf dem Smartphone) effektiv vor unbefugtem Zugriff.

\subsubsection{Mathematische Grundlage: Die Passwort-Entropie}
\label{subsec:entropie}
Die Stärke eines \gls{passwort}s gegenüber Angriffen lässt sich mathematisch über die Shannon-\gls{entropie} beschreiben. Sie gibt an, wie viele Versuche ein Angreifer im \mbox{schlechtesten Fall} benötigt, um den gesamten Suchraum zu durchsuchen.

Die Formel zur Berechnung der \gls{entropie} $E$ (in Bit) lautet:
\[
E = L \cdot \log_2(R)
\]

Dabei stehen die Variablen für folgende Werte:
\begin{itemize}
    \item $L$: Die Länge des \gls{passwort}s (Anzahl der verwendeten Zeichen).
    \item $R$: Die Grösse des verfügbaren Zeichensatzes (z.\,B. Kleinbuchstaben, Grossbuchstaben, Zahlen, Sonderzeichen).
\end{itemize}
\newpage
\subsubsection{Beispielrechnung}
\begin{itemize}
    \item \textbf{8 Zeichen (nur Kleinbuchstaben):} Bei $L = 8$ und einem Zeichensatz von $R = 26$ ergibt sich eine \gls{entropie} von:
    \[
    E_8 = 8 \cdot \log_2(26) \approx 37{,}6 \text{ Bit}
    \]
    \item \textbf{Verdopplung auf 16 Zeichen (gleicher Zeichensatz):} Bei $L = 16$ und $R = 26$ steigt die \gls{entropie} linear mit der Länge auf:
    \[
    E_{16} = 16 \cdot \log_2(26) \approx 75{,}2 \text{ Bit}
    \]
\end{itemize}

In \cref{tab:entropie_vergleich} sind die \gls{entropie}-Werte für verschiedene \gls{passwort}-Längen und Zeichensätze zusammenfassend dargestellt.

\begin{table}[htbp]
    \centering
    \begin{tabular}{|l|c|c|r|}
        \hline
        \textbf{Zeichensatz ($R$)} & \textbf{Länge ($L$)} & \textbf{Formel} & \textbf{\gls{entropie} ($E$)} \\
        \hline
        Kleinbuchstaben ($R=26$) & 8 & $8 \cdot \log_2(26)$ & $\approx 37{,}6$ Bit \\
        Kleinbuchstaben ($R=26$) & 16 & $16 \cdot \log_2(26)$ & $\approx 75{,}2$ Bit \\
        Alle Zeichen ($R=94$) & 8 & $8 \cdot \log_2(94)$ & $\approx 52{,}4$ Bit \\
        Kleinbuchstaben ($R=26$) & 20 & $20 \cdot \log_2(26)$ & $\approx 94{,}0$ Bit \\
        Alle Buchstaben ($R=52$) & 20 & $20 \cdot \log_2(52)$ & $\approx 114{,}0$ Bit \\
        Alle Buchstaben + Zahlen ($R=62$) & 20 & $20 \cdot \log_2(62)$ & $\approx 119{,}1$ Bit \\
        Alle Zeichen ($R=94$) & 20 & $20 \cdot \log_2(94)$ & $\approx 131{,}1$ Bit \\
        \hline
    \end{tabular}
    \caption{Vergleich der \gls{passwort}-\gls{entropie} für verschiedene Längen und Zeichensätze.}
    \label{tab:entropie_vergleich}
\end{table}

\subsubsection{Fazit der Berechnung}
Da die Anzahl der auszuprobierenden Kombinationen im Hintergrund exponentiell mit $2^E$ (E=\gls{entropie}) wächst, schützt das Verlängern des \gls{passwort}s um ein Vielfaches effektiver vor Angriffen als das krampfhafte Einfügen von Sonderzeichen bei einem zu kurzen \gls{passwort}. Zudem verhindert ein langes, zufälliges \gls{passwort}, dass ein \gls{dictionaryattack} erfolgreich ist.

\newpage

\subsubsection{Aktuelle Richtlinien und Best Practices}
Moderne Richtlinien haben sich von alten Regeln verabschiedet. Heute gelten folgende Hauptempfehlungen:
\begin{enumerate}
    \item \textbf{Länge vor Komplexität:} Bevorzugung von \gls{passphrase}n aus mehreren \mbox{zufälligen} Wörtern. Selbst wenn Angreifer raffinierte Wörterbuchangriffe oder Kombinationen ausprobieren, bricht ein System bei langen \gls{passphrase}n \mbox{(oft auch ohne erzwungene Sonderzeichen)} ein, weil die schiere Länge den Suchraum unknackbar macht (\gls{entropie} über 80 Bit).
    \item \textbf{Keine wiederholte Nutzung von \glslink{passwort}{Passwörtern}:} Jedes Konto sollte ein \mbox{einzigartiges} \gls{passwort} besitzen, um die Risiken bei \gls{datenleck}s zu minimieren.
    \item \textbf{Vermeidung von \gls{passwort}-Änderungen ohne Anlass:} Regelmässige \mbox{erzwungene} \gls{passwort}-Änderungen erhöhen die Wahrscheinlichkeit, dass Nutzer schwächere \glslink{passwort}{Passwörter} wählen oder diese aufschreiben. Nur bei Verdacht auf Kompromittierung sollte ein \gls{passwort} geändert werden.  
    \item \textbf{Einsatz von \gls{passwortmanager}:} Generierung und Verwaltung unikater \mbox{\gls{passwort}}.
    \item \textbf{Aktivierung von \gls{mfa}:} Einrichtung eines zweiten Faktors (z.\,B. aufs \mbox{Smartphone}).
\end{enumerate}

\newpage
% --- TEIL 2: PHISHING ---
\subsection{Phishing}
\label{sec:phishing}

\subsubsection{Einleitung}
In der digitalen Welt bilden nicht nur technische Barrieren wie \glslink{passwort}{Passwörter} die \mbox{Abwehrlinie}, sondern vor allem der Mensch selbst. Eines der erfolgreichsten und am \mbox{häufigsten} genutzten Einfallstore für Cyberkriminelle ist \gls{phishing}. Dabei wird gezielt die \mbox{menschliche Psyche} ausgenutzt, um an vertrauliche Daten zu gelangen oder \mbox{Schadsoftware} zu verbreiten \cite{bsi2023smtpsmuggling}.

\subsubsection{Grundlegende Begriffe und Angriffsarten}
\gls{phishing} gehört zu den klassischen \gls{socialengineering}-Angriffen. Statt komplexe \mbox{technische Systeme} direkt zu hacken, manipulieren Angreifer gezielt Personen:
\begin{itemize}
    \item \textbf{\gls{klassischesphishing}:} Täuschend echte E-Mails, die \mbox{vorgeben}, von \mbox{vertrauenswürdigen} \mbox{Institutionen} (wie Banken, Streaming-Diensten oder Behörden) zu \mbox{stammen}.
    \item \textbf{\gls{spearphishing}:} Ein gezielter Angriff (\gls{spearphishing}), der auf eine spezifische Person oder bestimmte Mitarbeiter zugeschnitten ist, um Vertrauen zu erwecken.
    \item \textbf{\gls{vishing} \& \gls{smishing}:} Betrugsversuche über Telefonanrufe (\textit{Voice \rightarrow \gls{vishing}}) oder \mbox{Kurznachrichten} (\textit{SMS \rightarrow \gls{smishing}}).
    \item \textbf{\gls{quishing}:} Das Platzieren von manipulierten QR-Codes (z. B. auf Plakaten oder Parkuhren), die zu gefälschten Webseiten führen.
\end{itemize}

\subsubsection{Psychologische Faktoren und Wirkungsweise}
Im Gegensatz zu mathematischen Angriffen zielt \gls{phishing} auf menschliche Schwachstellen wie Neugier, Hilfsbereitschaft, Gier oder vor allem \textbf{Angst und Zeitdruck} ab. Typische Maschen beinhalten drohende Konsequenzen wie angebliche Kontosperrungen, dringende Rechnungen oder verpasste Gewinne, damit das Opfer unüberlegt handelt.

\newpage
\subsubsection{Aktuelle Richtlinien und Best Practices}
Um sich vor \gls{phishing}-Angriffen zu schützen, gelten im Rahmen der IT-Sicherheit folgende zentrale Empfehlungen:

\begin{enumerate}
    \item \textbf{Kritische Hürden bewahren:} Niemals unüberlegt auf Links klicken oder \mbox{Anhänge} von unbekannten Absendern öffnen.
    \item \textbf{Absender und URLs genau prüfen:} Auf Unstimmigkeiten in E-Mail-Adressen (z. B. Buchstabendreher (\gls{typosquatting})) in Domains achten.
    \item \textbf{Multi-Faktor-Authentisierung (MFA) nutzen:} Selbst wenn Anmeldedaten \mbox{abgegriffen werden}, verhindert ein zusätzlicher Faktor (wie eine Bestätigung auf dem Smartphone) unbefugten Zugriff.
    \item \textbf{Keine sensiblen Daten preisgeben:} Seriöse Institutionen fragen niemals per E-Mail, SMS oder Telefon nach \glslink{passwort}{Passwörtern} oder \gls{pin}s.
    \item \textbf{Dringlichkeit und Druckmittel hinterfragen:} E-Mails, die mit \mbox{unbegründeten} Drohungen (z.~B. sofortige Kontosperrung) oder kurzen Fristen Panik erzeugen \mbox{wollen}, kritisch prüfen.
    \item \textbf{Sprache und Anrede analysieren:} Auf unpersönliche Anreden (z.~B. „Sehr \mbox{geehrter} Kunde“) sowie grammatikalische Fehler, fehlende Umlaute oder \mbox{ungewöhnliche} Formulierungen achten.
    \item \textbf{Direktzugriffe bevorzugen:} Webseiten und Portale stets manuell über die \mbox{Eingabe} der bekannten Adresse im Browser aufrufen, anstatt Links aus E-Mails zu folgen.
\end{enumerate}

\newpage
% --- TEIL 3: ÜBUNGSAUFGABEN ---
\subsection{Übungsaufgaben}
\label{sec:aufgaben}

Dieser Abschnitt enthält Vertiefungs- und Übungsaufgaben zu den Themen aus \mbox{\cref{sec:passwortsicherheit}} und \mbox{\cref{sec:phishing}}.

\subsubsection{Aufgabenblatt 1: Passwortsicherheit}
\label{subsec:aufgaben_passwort}

\subsubsubsection{Teil 1: Begriffsdefinitionen \& Grundlagen}

\begin{taskbox}[{Aufgabe 2.3.1.1.1: Begriffsdefinitionen}]
Erklären Sie kurz in Ihren eigenen Worten den Unterschied zwischen einem \newline \gls{bruteforce} und einem \gls{dictionaryattack}.
\end{taskbox}

\begin{solutionbox}[Musterlösung 2.3.1.1.1]
\begin{itemize}
    \item \textbf{\gls{bruteforce}:} Ein automatisierter Angriff, bei dem Programme \mbox{systematisch} alle möglichen Zeichenkombinationen durchprobieren.
    \item \textbf{\gls{dictionaryattack}:} Ein gezielter Angriff, bei dem Listen mit gebräuchlichen Wörtern, Namen und bekannten \glslink{passwort}{Passwörtern} \mbox{systematisch} getestet werden.
\end{itemize}
\end{solutionbox}

\begin{taskbox}[{Aufgabe 2.3.1.1.2: MFA}]
Was bedeutet die Abkürzung \gls{mfa} und wie schützt sie einen Account, selbst wenn ein Angreifer das \gls{passwort} herausgefunden hat?
\end{taskbox}

\begin{solutionbox}[Musterlösung 2.3.1.1.2]
\textbf{MFA} steht für \textbf{Multi-Faktor-Authentisierung}. Sie schützt einen Account, indem neben dem \gls{passwort} ein zweiter Faktor (z.\,B. eine Bestätigung auf dem Smartphone) erforderlich ist. Selbst wenn ein Angreifer das \gls{passwort} erlangt hat, bleibt der Zugriff ohne diesen zweiten Faktor verwehrt.
\end{solutionbox}

\newpage

\subsubsubsection{Teil 2: Mathematische Berechnungen (Entropie)}

Verwenden Sie die Formel zur Berechnung der \gls{entropie} aus \cref{subsec:entropie} \newline(sehen Sie auch \cref{tab:entropie_vergleich}):
\[
E = L \cdot \log_2(R)
\]

\begin{taskbox}[{Aufgabe 2.3.1.2.1: Einfluss des Zeichensatzes ($R$)}]
Gegeben ist ein \gls{passwort} der Länge $L = 10$.
\begin{enumerate}[label=\alph*)]
    \item Berechnen Sie die \gls{entropie} $E$, wenn nur Kleinbuchstaben verwendet werden \newline($R = 26$).
    \item Berechnen Sie die \gls{entropie} $E$, wenn Klein-, Grossbuchstaben, Zahlen und Sonderzeichen verwendet werden ($R = 94$).
\end{enumerate}
\end{taskbox}

\begin{solutionbox}[Musterlösung 2.3.1.2.1]
\begin{enumerate}[label=\alph*)]
    \item $E = 10 \cdot \log_2(26) \approx 10 \cdot 4{,}7004 = 47{,}0 \text{ Bit}$
    \item $E = 10 \cdot \log_2(94) \approx 10 \cdot 6{,}5546 = 65{,}5 \text{ Bit}$
\end{enumerate}
\end{solutionbox}

\begin{taskbox}[{Aufgabe 2.3.1.2.2: Länge vs. Komplexität}]
Betrachten Sie die folgenden beiden \glslink{passwort}{Passwörter}:
\begin{itemize}
    \item \textbf{Szenario A:} Ein komplexes \gls{passwort} mit 8 Zeichen ($L = 8$), \newline das \mbox{Kleinbuchstaben}, \mbox{Grossbuchstaben}, \mbox{Zahlen} und \mbox{Sonderzeichen} nutzt \newline($R = 94$).
    \item \textbf{Szenario B:} Ein zufälliges \gls{passwort} mit 20 Zeichen ($L = 20$), das nur aus Kleinbuchstaben besteht ($R = 26$).
\end{itemize}
\textbf{Aufgabe:} Berechnen Sie für beide Szenarien die \gls{entropie} $E$. Welches \gls{passwort} bietet eine höhere \gls{entropie} und warum entspricht dies der Regel „Länge vor \mbox{Komplexität}“?
\end{taskbox}

\begin{solutionbox}[Musterlösung 2.3.1.2.2]
\begin{itemize}
    \item \textbf{Szenario A ($L=8, R=94$):} $E_A = 8 \cdot \log_2(94) \approx 52{,}4 \text{ Bit}$
    \item \textbf{Szenario B ($L=20, R=26$):} $E_B = 20 \cdot \log_2(26) \approx 94{,}0 \text{ Bit}$
\end{itemize}
\textbf{Auswertung \& Begründung:} Szenario B bietet eine deutlich höhere \gls{entropie} ($94{,}0 \text{ Bit}$) als Szenario A ($52{,}4 \text{ Bit}$). Da der Suchraum im Hintergrund exponentiell mit $2^E$ (E=\gls{entropie}) wächst, vergrössert eine Zunahme der Länge ($L$) den Suchraum \mbox{wesentlich} stärker als ein grösserer Zeichensatz ($R$). Dies entspricht der Regel „Länge vor \mbox{Komplexität}“.
\end{solutionbox}

\newpage

\subsubsubsection{Teil 3: Anwendung \& Best Practices}

\begin{taskbox}[{Aufgabe 2.3.1.3.1: Richtlinien beurteilen}]
Welche der folgenden Aussagen entsprechen den aktuellen Best Practices \newline(z.\,B. nach \gls{nist})? Begründen Sie Ihre Antwort kurz.
\begin{enumerate}[label=\Alph*)]
    \item „Ein sicheres \gls{passwort} muss alle 30 Tage geändert werden.“
    \item „\gls{passphrase}n aus mehreren zufälligen Wörtern sind in der Regel sicherer und merkbarer als kurze, komplexe Zeichenketten.“
    \item „Verwenden Sie das gleiche \gls{passwort} für mehrere Dienste, solange es sehr lang ist.“
    \item „Der Einsatz eines \gls{passwortmanager} wird empfohlen, um für jeden Dienst ein unikat zufälliges \gls{passwort} zu generieren.“
\end{enumerate}
\end{taskbox}

\begin{solutionbox}[Musterlösung 2.3.1.3.1]
\begin{enumerate}[label=\Alph*)]
    \item \textbf{Falsch:} Regelmässige erzwungene Änderungen verleiten Nutzer zu \mbox{schwächeren} \glslink{passwort}{Passwörtern} oder Mustern. Eine Änderung sollte nur bei Verdacht auf Kompromittierung erfolgen.
    \item \textbf{Richtig (Best Practice):} \gls{passphrase}n bieten durch ihre Länge eine hohe \gls{entropie} (über 80 Bit) und sind widerstandsfähiger gegen Angriffe.
    \item \textbf{Falsch:} \glslink{passwort}{Passwörter} dürfen niemals mehrfach genutzt werden, um \mbox{Folgeschäden} bei Datenlecks zu verhindern.
    \item \textbf{Richtig (Best Practice):} Ein \gls{passwortmanager} ermöglicht die sichere \mbox{Generierung} und Verwaltung von einzigartigen \glslink{passwort}{Passwörtern} für jeden Dienst.
\end{enumerate}
\end{solutionbox}

\newpage

\subsubsection{Aufgabenblatt 2: Phishing, Social Engineering \& Best Practices}
\label{subsec:aufgaben_phishing}

\subsubsubsection{Teil 1: Verständnis \& Begriffe}

\begin{taskbox}[{Aufgabe 2.3.2.1.1: Begriffsdefinitionen \& Abgrenzung}]
Erklären Sie kurz in Ihren eigenen Worten den Unterschied zwischen \glslink{klassischesphishing}{klassischem Phishing} und gezieltem \gls{spearphishing}. Gehen Sie dabei darauf ein, wer jeweils das Ziel des Angriffs ist.
\end{taskbox}

\begin{solutionbox}[Musterlösung 2.3.2.1.1]
\begin{itemize}
    \item \textbf{\gls{klassischesphishing}:} Richtet sich ungezielt an eine breite Masse von \mbox{Personen} \newline (z.\,B. gefälschte Massen-E-Mails im Namen bekannter Unternehmen).
    \item \textbf{\gls{spearphishing}:} Ein hochspezifischer, massgeschneiderter Angriff auf eine einzelne ausgewählte Person oder eine konkrete Mitarbeitergruppe innerhalb einer Organisation.
\end{itemize}
\end{solutionbox}

\begin{taskbox}[{Aufgabe 2.3.2.1.2: Gefahrenkanäle (Vishing, Smishing \& Quishing)}]
Cyberkriminelle nutzen nicht nur E-Mails. Ordnen Sie den folgenden Angriffsformen die passende Methode zu und erklären Sie kurz, worauf man achten muss:
\begin{enumerate}[label=\alph*)]
    \item \textbf{\gls{vishing}}
    \item \textbf{\gls{smishing}}
    \item \textbf{\gls{quishing}}
\end{enumerate}
\end{taskbox}

\begin{solutionbox}[Musterlösung 2.3.2.1.2]
\begin{enumerate}[label=\alph*)]
    \item \textbf{\gls{vishing} (Voice Phishing):} Phishing über Telefonanrufe. \newline Beachtung: Seriöse Anrufer verlangen niemals \glslink{passwort}{Passwörter} oder \gls{pin}s am \mbox{Telefon.}
    \item \textbf{\gls{smishing} (SMS Phishing):} Betrug über Kurznachrichten (SMS). \newline Beachtung: Keine unüberlegten Klicks auf SMS-Links und keine Preisgabe vertraulicher Daten.
    \item \textbf{\gls{quishing} (QR-Code Phishing):} Platzierung gefälschter QR-Codes. \newline Beachtung: Ziel-URL vor dem Aufrufen im Scanner genau prüfen.
\end{enumerate}
\end{solutionbox}

\newpage

\begin{taskbox}[{Aufgabe 2.3.2.1.3: Multi-Faktor-Authentisierung (MFA)}]
Warum ist \gls{mfa} eines der effektivsten Schutzmittel gegen Phishing? \newline Wie schützt \gls{mfa} einen Account, selbst wenn ein Opfer seine Anmeldedaten auf einer gefälschten Phishing-Website eingegeben hat?
\end{taskbox}

\begin{solutionbox}[Musterlösung 2.3.2.1.3]
MFA ist extrem effektiv, da \glslink{passwort}{Passwörter} alleine für eine erfolgreiche Übernahme nicht ausreichen. Selbst wenn ein Opfer sein \gls{passwort} auf einer Phishing-Webseite eingibt, scheitert der Angreifer beim Anmeldeversuch, weil die Bestätigung über den zweiten Faktor (z.\,B. Smartphone-App) fehlt.
\end{solutionbox}

\newpage

\subsubsubsection{Teil 2: Psychologische Faktoren \& Szenario-Analyse}

\begin{taskbox}[{Aufgabe 2.3.2.2.1: Psychologische Tricks}]
\gls{socialengineering} nutzt primär menschliche Schwachstellen statt technischer \mbox{Systemlücken}. Nennen Sie drei psychologische Faktoren, die \mbox{Angreifer} \mbox{gezielt} \mbox{ausnutzen,} und beschreiben Sie kurz, wie diese einen Menschen zu \mbox{unüberlegtem} \mbox{Handeln} verleiten.
\end{taskbox}

\begin{solutionbox}[Musterlösung 2.3.2.2.1]
\begin{enumerate}
    \item \textbf{Angst:} Drohungen wie Kontosperrungen erzeugen Panik, sodass Logik und Vorsicht ausgeschaltet werden.
    \item \textbf{Zeitdruck / Dringlichkeit:} Kurze Fristen verleiten dazu, übereilt und \mbox{unüberlegt} zu handeln.
    \item \textbf{Neugier / Gier:} Versprechen von Gewinnen oder exklusiven Informationen verleiten zum unkritischen Klicken auf Links.
\end{enumerate}
\end{solutionbox}

In \cref{fig:phishing_email} ist eine typische verdächtige Nachricht für die nachfolgende \newline Szenario-Analyse dargestellt.

\begin{figure}[H]
\centering
\begin{emailbox}[Phishing-E-Mail-Verdachtsfall]
\textbf{Von:} \texttt{support-service@bank-security-check.ch} \\
\textbf{An:} \texttt{max.muster@kanti.ch} \\
\textbf{Betreff:} \textbf{DRINGEND: Ihr Konto wird in 2 Stunden gesperrt!} \\[0.5em]
Sehr geehrter Kunde, \\[0.5em]
aufgrund verdächtiger Aktivitäten müssen Sie Ihre Daten sofort verifizieren. \newline Wenn Sie nicht innerhalb von 120 Minuten handeln, wird Ihr Zugang dauerhaft gesperrt. \\[0.5em]
Klicken Sie hier: \url{http://www.kanton-bank-login-sicherheit.com/verify}
\end{emailbox}
\caption{Verdächtiges E-Mail-Fallbeispiel für die Szenario-Analyse.}
\label{fig:phishing_email}
\end{figure}

\newpage

\begin{taskbox}[{Aufgabe 2.3.2.2.2: Fallbeispiel – Szenario-Analyse}]
Lesen Sie die in \cref{fig:phishing_email} dargestellte Nachricht aufmerksam durch.
\begin{enumerate}[label=\alph*)]
    \item Identifizieren Sie mindestens drei Merkmale in dieser Nachricht, die eindeutig auf einen Phishing-Angriff hinweisen.
    \item Welche Form von Phishing und welcher psychologische Mechanismus liegen hier vor?
\end{enumerate}
\end{taskbox}

\begin{solutionbox}[Musterlösung 2.3.2.2.2]
\begin{enumerate}[label=\alph*)]
    \item \textbf{Drei Verdachtsmerkmale:}
    \begin{itemize}
        \item \textit{Verdächtige Absender- und Ziel-URL:} \newline Weder \texttt{@bank-security-check.ch} noch die Domain \texttt{kanton-bank-login-sicherheit.com} stammen von einer offiziellen Bank.
        \item \textit{Unpersönliche Anrede:} Generische Anrede „Sehr geehrter Kunde“.
        \item \textit{Massiver Zeitdruck und Drohung:} Fristsetzung von 120 Minuten mit \mbox{Androhung} einer dauerhaften Kontosperrung.
    \end{itemize}
    \item \textbf{Form \& Psychologie:} Es handelt sich um \textbf{klassisches Phishing}. \newline Der psychologische Mechanismus basiert auf Angst und künstlich erzeugtem Zeitdruck.
\end{enumerate}
\end{solutionbox}

\newpage

\subsubsubsection{Teil 3: Anwendung \& Best Practices (Richtlinien beurteilen)}

\begin{taskbox}[{Aufgabe 2.3.2.3.1: Richtlinien beurteilen}]
Beurteilen Sie die folgenden Aussagen. Ist die Aussage \textbf{Richtig} oder \textbf{Falsch}? \newline Begründen Sie Ihre Antwort jeweils kurz in einem bis drei Sätzen.

\textbf{Aussage A:} „Seriöse Banken und Institutionen fordern Kunden gelegentlich per E-Mail auf, ihr \gls{passwort} einzugeben, um die Identität bei Sicherheitsupdates zu bestätigen.“

\textbf{Aussage B:} „Beim Aufruf von Links sollte stets die exakte URL (Internetadresse) geprüft werden, da Angreifer häufig Buchstabendreher oder gefälschte Subdomains (\gls{typosquatting}) nutzen.“

\textbf{Aussage C:} „Ein Phishing-Angriff kann nur dann erfolgreich sein, wenn auf einem Computer keine Antiviren-Software installiert ist.“

\textbf{Aussage D:} „Das Scannen von QR-Codes auf Plakaten oder Parkuhren ist völlig risikofrei, da QR-Codes keine Viren oder Links enthalten können.“
\end{taskbox}

\begin{solutionbox}[Musterlösung 2.3.6.1]
\begin{itemize}
    \item \textbf{Aussage A: Falsch.} Seriöse Institutionen fragen niemals per E-Mail, SMS oder Telefon nach \glslink{passwort}{Passwörtern} oder \gls{pin}s.
    \item \textbf{Aussage B: Richtig.} Angreifer versuchen oft, Opfer durch täuschend \mbox{ähnliche} Domains oder gefälschte Subdomains (Typosquatting) in die Irre zu \mbox{führen.}
    \item \textbf{Aussage C: Falsch.} Phishing nutzt primär die menschliche Psyche (\gls{socialengineering}). Eine Antiviren-Software schützt nicht davor, dass der Nutzer seine Daten selbst auf einer gefälschten Website eingibt.
    \item \textbf{Aussage D: Falsch.} Beim \gls{quishing} werden QR-Codes auf schädliche oder gefälschte Zielseiten verlinkt, was ein erhebliches Sicherheitsrisiko darstellt.
\end{itemize}
\end{solutionbox}

\end{document}